% ============================================================
%  SCADA Hub — Manager Punchlist (ADO Submission)
% ============================================================
\documentclass[10pt]{article}
\usepackage[T1]{fontenc}
\usepackage[utf8]{inputenc}
\usepackage{lmodern}
\usepackage{microtype}
\usepackage{palatino}
\usepackage[scaled=0.82]{beramono}
\usepackage[top=0.55in, bottom=0.50in, left=0.65in, right=0.65in]{geometry}
\usepackage[dvipsnames,svgnames,x11names]{xcolor}
\usepackage{tikz}
\usepackage{enumitem}
\usepackage[most]{tcolorbox}
\usepackage{multicol}
\usepackage[colorlinks=true, linkcolor=sblue, urlcolor=scyan,
  pdftitle={SCADA Hub --- Manager Punchlist}]{hyperref}
\usepackage{parskip}

\definecolor{sblue}{HTML}{2563EB}
\definecolor{scyan}{HTML}{3B82F6}
\definecolor{snavy}{HTML}{1E3A5F}
\definecolor{sorange}{HTML}{EA580C}
\definecolor{sgreen}{HTML}{15803D}
\definecolor{sred}{HTML}{DC2626}
\definecolor{samber}{HTML}{D97706}
\definecolor{sgray}{HTML}{6B7280}
\definecolor{slightgray}{HTML}{F1F5F9}

\setlength{\parskip}{2pt}
\setlength{\parindent}{0pt}
\pagestyle{empty}

% Priority badge helper
\newcommand{\high}{\textcolor{sred}{\small\bfseries HIGH}}
\newcommand{\med}{\textcolor{samber}{\small\bfseries MED}}
\newcommand{\low}{\textcolor{sgreen}{\small\bfseries LOW}}
\newcommand{\done}{\textcolor{sgreen}{\small\bfseries DONE}}

\begin{document}

%% ---- HEADER ----
\begin{tikzpicture}[remember picture, overlay]
  \fill[snavy] (current page.north west) rectangle ([yshift=-1.0in]current page.north east);
  \node[anchor=west, text=white, font=\fontsize{20}{24}\bfseries\sffamily]
    at ([xshift=0.65in, yshift=-0.35in]current page.north west) {SCADA Hub};
  \node[anchor=west, text=scyan!80, font=\fontsize{11}{13}\sffamily]
    at ([xshift=0.65in, yshift=-0.63in]current page.north west)
    {Work Item Punchlist · ADO Submission · 2026-05-18};
  \draw[sorange, line width=2pt]
    ([xshift=0.65in, yshift=-0.83in]current page.north west) --
    ([xshift=3.5in,  yshift=-0.83in]current page.north west);
\end{tikzpicture}

\vspace{0.70in}

{\small\sffamily\color{sgray}\MakeUppercase{Overview}}
\smallskip

{\small 8 study guide sites covering Modbus, OPC UA, DNP3, IEC 61131-3, PID Controllers, SEL RTAC, Ignition SCADA, and Wireshark. Each guide is a React/Vite SPA deployed to GitHub Pages, with quizzes, flashcards, a Code Lab, and a downloadable course PDF. The hub landing page (scada-hub.vercel.app) links all 8 guides.}

\smallskip
{\color{snavy}\rule{\linewidth}{1.2pt}}
\smallskip

%% ---- SECTION: FEATURES ----
{\small\sffamily\color{sgray}\MakeUppercase{Feature Work}}
\smallskip

\setlist[enumerate]{noitemsep, topsep=2pt, leftmargin=*, label=\textcolor{sblue}{\small\bfseries\arabic*.}}
\begin{enumerate}\small
  \item \textbf{Quiz Question Hints} \hfill \high \\
    Add a \texttt{hint} field to every quiz question across all 8 guides' \texttt{quizzes.js} data files — each hint points to WHERE the answer lives (chapter name, section heading, flashcard ID). Wire up a "Hint" button in \texttt{Quiz.jsx} that reveals the source on click. Touches all 8 guides × 600–900 questions each.

  \item \textbf{Certification Relevance Sidebar Panel} \hfill \med \\
    Build a new sidebar button per guide (modal or slide-out panel) with the cert details: name, exam body, link, description, and how each course chapter maps to exam domains. This content got cut from the syllabus PDF — the sidebar panel is its permanent home. Touches all 8 guides.

  \item \textbf{DNP3 Study Guide LaTeX PDF} \hfill \med \\
    Write, compile, and deploy the DNP3 full-length study guide PDF (same scope as the Modbus PDF). Drop it at \texttt{site/public/study\_guide.pdf}. Right now the sidebar PDF button returns a 404. DNP3 guide only.

  \item \textbf{Wrong-Answer Tracking + Review Missed Mode} \hfill \med \\
    Track which question IDs a user missed per chapter — Supabase is already wired for this. Add a "Review Missed" quiz mode that serves only the questions they got wrong. Touches all 8 guides + Supabase schema.

  \item \textbf{Spaced Repetition on Flashcards} \hfill \low \\
    Swap out the current random shuffle for the SM-2 algorithm (\textasciitilde150 lines of JS). Missed cards come back sooner; cards a user has nailed surface less often. Touches all 8 guides' \texttt{Flashcards.jsx}.

  \item \textbf{Timed Exam Mode} \hfill \low \\
    Hub-level feature: 60 questions, 90-minute timer, pulled randomly across all guides in a topic area (OT protocols, ICS security, etc.). Supabase handles scoring and the leaderboard. Touches scada-hub + Supabase.

  \item \textbf{Cross-Guide "See Also" Callouts} \hfill \low \\
    Add "See Also" callouts inside chapters that link out to related chapters in other guides (e.g., Wireshark Modbus chapter → Modbus guide, OPC UA Security → Wireshark OPC UA chapter). Touches all 8 guides' content pages.
\end{enumerate}

\smallskip
{\color{snavy}\rule{\linewidth}{0.4pt}}
\smallskip

%% ---- SECTION: CONTENT GAPS ----
{\small\sffamily\color{sgray}\MakeUppercase{Content Gaps — Syllabus vs.\ Coursework}}
\smallskip

\begin{enumerate}[resume]\small
  \item \textbf{Modbus: Prerequisites + ModRSsim2 + TCP Flashcards} \hfill \med \\
    (a) Add a prerequisite statement to \texttt{Home.jsx} / Intro chapter. (b) Call out ModRSsim2 in the TCP or Required Software section of the web guide. (c) Beef up the TCP chapter flashcards — currently 3; the syllabus implies 5+, covering transaction ID pipelining and Unit ID in gateways.

  \item \textbf{OPC UA: Sessions/Channels Chapter + Discovery Section} \hfill \med \\
    (a) Add a Sessions \& Channels section to \texttt{Architecture.jsx} — token renewal, reconnection retry, ActivateSession flow. (b) Add a Discovery section: GetEndpoints, FindServers, GDS pattern, certificate trust flow. (c) Get open62541 into at least one chapter page — right now it only shows up in quiz questions.

  \item \textbf{IEC 61131-3: PLCopen FB Library + Vendor Portability} \hfill \med \\
    (a) Add a section on PLCopen standard FB signatures (TON/TOF/CTU), instance management, and how to navigate a library. (b) Add a vendor portability section on what breaks when you move code across Siemens/Beckhoff/SEL. (c) Add a table showing which IEC version applies to which platforms.

  \item \textbf{PID: Scilab Lab Exercises + Laplace Prerequisites} \hfill \med \\
    (a) Add Scilab step-test simulation exercises to the Process Characterization chapter. (b) Drop a note in the Digital Implementation chapter flagging the Laplace/transfer-function prerequisites — students without calculus need to know what to skip or review before they hit that section.

  \item \textbf{DNP3: Prerequisites Statement in Web Guide} \hfill \low \\
    Add a short prerequisite note to \texttt{Home.jsx} or the Intro chapter: basic TCP/IP and Modbus familiarity assumed. All 10 topics are fully covered — everything else checks out.

  \item \textbf{RTAC: Modbus Gateway Chapter} \hfill \med \\
    Add a dedicated Modbus Gateway chapter (or at minimum a major section) covering DNP3-to-Modbus point mapping, register layout design, data type conversion rules, and how translation failures present. Syllabus Topic 4 ("Modbus Gateway") has no matching web chapter right now. Also flesh out I/O module selection and pull the SEL-5033 virtual RTAC into the lab exercise.

  \item \textbf{Ignition: OPC-UA Integration Chapter Depth} \hfill \med \\
    Audit \texttt{OpcUa.jsx} against syllabus Topic 3 (OPC-UA Integration). Check that the PLC-to-live-tags workflow and Ignition-as-OPC-UA-server config are fully covered. Add a SQL prerequisite note to Home.jsx. (\emph{Audit pending.})

  \item \textbf{Wireshark: SSH Remote Capture + tcpreplay} \hfill \low \\
    Syllabus Topic 9 (Advanced Techniques) calls out SSH remote capture and tcpreplay. Verify both are in \texttt{Advanced.jsx} and add anything that's missing. (\emph{Audit pending.})
\end{enumerate}

\smallskip
{\color{snavy}\rule{\linewidth}{0.4pt}}
\smallskip

%% ---- SECTION: BUGS + INFRASTRUCTURE ----
{\small\sffamily\color{sgray}\MakeUppercase{Bugs \& Infrastructure}}
\smallskip

\begin{enumerate}[resume]\small
  \item \textbf{Auth: Duplicate Email Signup Error} \hfill \med \\
    Registering with an existing email drops a raw Supabase error on the user. Show a plain "account already exists" message with a "Reset password" link, and wire up the reset-password flow from that error state. Scope: scada-hub \texttt{AuthPage.jsx}.

  \item \textbf{Quiz Reports RLS — Supabase} \hfill \med \\
    Right now any logged-in user can read any other user's quiz history. Add RLS policies to the \texttt{quiz\_results} table: users can only see their own rows; admin gets a bypass policy.

  \item \textbf{Mobile UX Audit} \hfill \med \\
    Mobile layout is rough across all 8 guides and the hub. Needs a full pass: text sizes, tap targets, sidebar behavior, LearningPath horizontal scroll on small screens.

  \item \textbf{Supabase Email Confirmation Redirect} \hfill \low \\
    The confirmation email link sends users to a broken URL. Fix: update Site URL and Redirect URLs in Supabase Auth settings to point to \texttt{scada-hub.vercel.app}.

  \item \textbf{LearningPath Hover Scale Bug} \hfill \low \\
    Hovering a step card fires \texttt{whileHover=\{scale: 1.03\}} — the card pops out and text gets hard to read. Dial back or remove the scale transform. Scope: \texttt{scada-hub/site/src/components/LearningPath.jsx}.

  \item \textbf{RTAC PDF Layout Overhaul} \hfill \low \\
    The RTAC study guide PDF has bad spacing, orphaned headings, and an intro style that doesn't match the Modbus PDF. Needs a formatting pass to bring it in line with the Modbus LaTeX template.

  \item \textbf{OG Image Quality Improvement} \hfill \low \\
    OG images are generated with Python PIL and the system Helvetica — they work but look rough. Replace with a custom SVG or Vercel's \texttt{@vercel/og} for server-side generation at deploy time. Touches scada-hub \texttt{public/og-image.png} + all 8 guide repos.
\end{enumerate}

\smallskip
{\color{snavy}\rule{\linewidth}{0.4pt}}
\smallskip

{\small\sffamily\color{sgray}\MakeUppercase{UI / Branding}}
\smallskip

\begin{enumerate}[resume]\small
  \item \textbf{Login Page: Company A Brand Colors} \hfill \med \\
    The SCADA Hub login/auth page (\texttt{ScadaBackground.jsx} + \texttt{AuthPage.jsx}) uses orange as its primary accent. Swap it for whatever hex Company A provides. Touches \texttt{scada-hub/site/src/components/ScadaBackground.jsx} and \texttt{AuthPage.jsx}. Get the target hex before starting.

  \item \textbf{Login Page: Floating Telemetry Values} \hfill \med \\
    On the animated SCADA login background, the live telemetry tickers (PV kW, SOC\%, RTT ms, MVA values) overlap node icons on some screen sizes. Fix: (a) reposition ticker boxes so they don't sit on top of the icon, and (b) add a slow vertical float animation to give the background a bit more life. Touches \texttt{ScadaBackground.jsx} node layout and ticker positioning logic.

  \item \textbf{All Courses Sidebar Panel — Design Polish} \hfill \low \\
    The "All Courses" collapsible nav item landed in all 8 guide sidebars with guide icons and Vercel URLs. Follow-up polish: (a) highlight the current guide in the dropdown so users can tell where they are, (b) consider showing a chapter count or question count as a subtitle under each guide name.

  \item \textbf{Navigation: Sticky Prev/Next Footer + Keyboard Shortcuts} \hfill \high \\
    Prev/Next chapter navigation exists only at the bottom of each page. Best practice (GitBook, ReadTheDocs, Stripe Docs): add a sticky footer bar visible at all times showing the current chapter name and Prev/Next buttons. Add keyboard shortcuts: left/right arrow for prev/next chapter. This is the single highest-impact navigation improvement — users on long chapters have to scroll to the bottom to move forward. Touches \texttt{ChapterLayout.jsx} across all 8 guides.

  \item \textbf{Navigation: Chapter Completion UX} \hfill \high \\
    There is no explicit "mark complete" UI. Reading a chapter currently auto-marks it as visited via \texttt{markChapterVisited()} on mount — but users have no way to know this happened and no way to un-mark. Best practice: add a "Mark Complete" button at the bottom of each chapter (or auto-trigger on scroll-to-end) with visible confirmation (check animation, chapter turns green in sidebar). Also: when a user completes a chapter, highlight the next chapter in the sidebar and surface a "Next Up" card at the bottom of the page above the nav arrows. Touches all 8 guides' \texttt{ChapterLayout.jsx} and \texttt{Sidebar.jsx}.

  \item \textbf{Navigation: Sidebar Progress Dots — Clarity} \hfill \med \\
    The three dots per chapter (quiz L1/L2/L3 passed) are subtle and the meaning is non-obvious without a legend. Best practice: replace or supplement with a checkmark icon for fully-visited chapters (quiz optional) and add a one-line legend below the progress bar. The visited state (amber dot) and passed states (green/amber/rose dots) are not clearly differentiated. Research shows checkmarks are universally understood; colored dots require learning. Touches \texttt{Sidebar.jsx} across all 8 guides.

  \item \textbf{Navigation: Mobile Bottom Tab Bar Exploration} \hfill \med \\
    Current mobile: hamburger button top-left, full-screen backdrop drawer. Apple/Google HIG and NNG both recommend bottom tab bars for primary navigation in learning apps — thumb-reachable, persistent, no hidden state. Explore replacing the hamburger with a bottom tab bar for: Home, Chapters (opens drawer), Flashcards, Quiz, Progress. This would be a significant UX improvement for mobile users who represent a growing share. Touches \texttt{Sidebar.jsx} and \texttt{App.jsx} across all 8 guides.

  \item \textbf{Navigation: In-Page Table of Contents} \hfill \low \\
    Longer chapters (DataTypes, FunctionCodes, Security, Troubleshoot) have 5–8 sections without in-page anchors. Best practice for long-form content: add a sticky in-page TOC that highlights the current section on scroll. GitBook and Stripe Docs both use this. Implementation: generate TOC from \texttt{h2} elements, scroll-position listener, sticky right panel or top-of-page collapsible. Only worth doing for chapters that exceed 4 major sections. Touches \texttt{ChapterLayout.jsx} or a new \texttt{TableOfContents.jsx} component.
\end{enumerate}

\smallskip
{\color{snavy}\rule{\linewidth}{0.4pt}}
\smallskip

{\small\sffamily\color{sgray}\MakeUppercase{Course Audit — Modbus Gaps}}
\smallskip

\begin{enumerate}[resume]\small
  \item \textbf{Modbus Syllabus: Add RS-485 as Standalone Topic} \hfill \low \\
    The syllabus lists 11 topics but the web guide and LaTeX PDF both treat RS-485 Physical Layer as a full standalone chapter (Ch 4). Add RS-485 as Topic 3 in the syllabus (between Data Model and RTU Mode) — brings the total to 12 topics and matches the actual chapter count.

  \item \textbf{Modbus Web: Add FC08 and FC43 to Function Codes Chapter} \hfill \med \\
    \texttt{FunctionCodes.jsx} only covers FC01, FC02, FC03, FC04, FC05, FC06, FC15, FC16. The LaTeX PDF (Ch 8) and flashcards both include FC08 (Diagnostics) and FC43 (Encapsulated Interface Transport / MEI). Add those two to the web chapter so all three sources stay in sync.

  \item \textbf{Modbus Web: Add Exception Codes 10 and 11 to Exceptions Chapter} \hfill \med \\
    \texttt{Exceptions.jsx} stops at codes 01--08. The LaTeX PDF (Ch 9) and flashcards both cover gateway codes 10 (Gateway Path Unavailable) and 11 (Gateway Target Device Failed to Respond). Add them to the web chapter and fix the syllabus wording — it currently says ``eight exception identifiers'' but there are eleven.

  \item \textbf{Modbus Web: Add Connection Management to TCP Chapter} \hfill \low \\
    \texttt{TCP.jsx} doesn't cover stale connections or multiple simultaneous connections. LaTeX Ch 7 covers both. Add a Connection Management section to the web chapter.

  \item \textbf{Modbus Web: Add Poll Rate Overrun to Troubleshooting Chapter} \hfill \low \\
    The syllabus explicitly calls out ``poll rate overruns'' as a troubleshooting topic but it's nowhere in \texttt{Troubleshoot.jsx}. Add it as a symptom scenario: device busy exceptions, missed polls, sluggish HMI response — cause is too many registers per cycle or too short a poll interval for the baud rate.

  \item \textbf{Modbus Web Lab: Add Gateway Exception 11 and Security Audit Scenarios} \hfill \med \\
    LaTeX Ch 12 has Scenario 5 (Modbus TCP Gateway Exception 11) and Scenario 6 (Security Audit Finding) — neither is in \texttt{Lab.jsx}. Add both as interactive exercises or frame-decode problems to close the gap.

  \item \textbf{Modbus Flashcards: Expand Security Chapter} \hfill \low \\
    The Ch 11 Security flashcard set has only 4 cards even though the chapter covers 5 attack vectors, 5 defense-in-depth layers, and 4 NERC CIP standards. Shoot for 8--10 cards to match the depth of other chapters.
\end{enumerate}

\smallskip
{\color{snavy}\rule{\linewidth}{0.4pt}}
\smallskip

{\small\sffamily\color{sgray}\MakeUppercase{Course Audit — OPC UA Gaps}}
\smallskip

\begin{enumerate}[resume]\small
  \item \textbf{OPC UA Web: Add Dedicated Discovery Chapter} \hfill \med \\
    The syllabus has Discovery as Topic 8 — Local Discovery Server, Global Discovery Server, FindServers, GetEndpoints, endpoint enumeration, open-discovery security risk. No \texttt{Discovery.jsx} page exists; the topic shows up briefly in \texttt{Intro.jsx} and in one flashcard (t10). Build a full web chapter covering LDS/GDS architecture, FindServers vs GetEndpoints, certificate trust during discovery, and the security risk of unauthenticated endpoint enumeration.

  \item \textbf{OPC UA Flashcards: Add Troubleshooting Cards} \hfill \med \\
    \texttt{Troubleshoot.jsx} is a solid chapter — 7 status codes, a 5-step diagnostic ladder, UA Expert workflow, Wireshark analysis — but zero flashcards cover any of it. Add 6--8 cards hitting the most common status codes (\texttt{BadCertificateUntrusted}, \texttt{BadNodeIdUnknown}, \texttt{BadUserAccessDenied}, etc.) and the diagnostic ladder sequence.

  \item \textbf{OPC UA Flashcards: Add Ignition Integration Cards} \hfill \low \\
    \texttt{IgnitionIntegration.jsx} covers Ignition port 62541, the 5-step certificate trust workflow, and SEL RTAC as an OPC UA server — none of it appears in flashcards. Add 4--5 cards on Ignition-specific OPC UA behavior.

  \item \textbf{OPC UA Web Security Chapter: Add CRL and Role-Based Security} \hfill \low \\
    Flashcards sec05 (Certificate Revocation Lists) and sec07 (role-based access: Observer, Operator, Engineer, Supervisor) exist but neither topic shows up in \texttt{Security.jsx}. Add subsections for CRL checking and built-in OPC UA roles to bring the web chapter in line with the flashcards.

  \item \textbf{OPC UA Web Services Chapter: Add RegisterNodes} \hfill \low \\
    \texttt{RegisterNodes}/\texttt{UnregisterNodes} is a performance optimization service that appears in flashcard s08 and Level 3 quiz questions but is absent from \texttt{Services.jsx}. Add a short section on when and why to use \texttt{RegisterNodes} for high-frequency polling.
\end{enumerate}

\smallskip
{\color{snavy}\rule{\linewidth}{0.4pt}}
\smallskip

{\small\sffamily\color{sgray}\MakeUppercase{Course Audit — DNP3 Gaps}}
\smallskip

\begin{enumerate}[resume]\small
  \item \textbf{DNP3 Flashcards: Add Application Layer Chapter Cards} \hfill \med \\
    The \texttt{appLayer} chapter has a full quiz question set but zero flashcards. Add 5--7 cards covering the DNP3 Application Layer: request/response fragments, function codes (READ, WRITE, SELECT, OPERATE, DIRECT OPERATE), IIN bits, sequence numbering, and fragmentation reassembly. That's the only gap in the DNP3 audit — all 10 syllabus topics are otherwise covered.
\end{enumerate}

\smallskip
{\color{snavy}\rule{\linewidth}{0.4pt}}
\smallskip

{\small\sffamily\color{sgray}\MakeUppercase{Course Audit — IEC 61131-3 Gaps}}
\smallskip

\begin{enumerate}[resume]\small
  \item \textbf{IEC 61131-3 Flashcards: Add Lab Chapter Cards} \hfill \med \\
    The Lab \& Practice chapter (\texttt{Ch 10}) has 60 quiz questions and zero flashcards. Add 6--8 cards covering the practical lab scenarios: PLC program structure, scan cycle, I/O addressing syntax, and the beginner mistakes that show up most often in the quiz scenario questions.

  \item \textbf{IEC 61131-3 Flashcards: Split Languages Category by Language} \hfill \med \\
    The \texttt{languages} flashcard category has \textasciitilde47 cards crammed across four distinct chapters (Intro, Structured Text, Ladder Diagram, Function Block Diagram) — each of those chapters has 60 quiz questions but shares one thin flashcard set. Split into at least two categories: \texttt{ld-fbd} (graphical languages) and \texttt{st-il} (text languages), and expand to \textasciitilde20 cards per language to match quiz depth.

  \item \textbf{IEC 61131-3 Flashcards: Cover Scenario-Type Questions} \hfill \low \\
    The quiz data has \textasciitilde80 scenario questions (\texttt{iec\_*\_sc} entries, \textasciitilde5 per chapter) covering applied troubleshooting and design decisions. No flashcard category touches scenario-style content. Add a \texttt{scenarios} category with 10--12 cards framing the most instructive setups and how to work through them.
\end{enumerate}

\smallskip
{\color{snavy}\rule{\linewidth}{0.4pt}}
\smallskip

{\small\sffamily\color{sgray}\MakeUppercase{Course Audit — PID Controllers Gaps}}
\smallskip

\begin{enumerate}[resume]\small
  \item \textbf{PID Flashcards: Add Cards for 5 Uncovered Chapters} \hfill \high \\
    5 of 10 quiz chapters have zero flashcards: Closed-Loop Intro (65 Qs), Loop Fundamentals (65 Qs), P/I/D Action (65 Qs), Digital PID (65 Qs), and Troubleshooting (65 Qs). That's roughly half the quiz content with nothing to reinforce it. Priority order: (1) P/I/D Action, (2) Troubleshooting, (3) Digital PID, (4) Intro, (5) Loop Fundamentals. Shoot for 6--8 cards per chapter.

  \item \textbf{PID Flashcards: Align Category Names to Chapter Boundaries} \hfill \low \\
    The current flashcard categories (fundamentals, tuning, dynamics, advanced, practical) don't map to the 10 syllabus chapters. A student coming off a quiz chapter can't find the matching flashcard deck. Rename or split categories so each quiz chapter has a clear flashcard counterpart.

  \item \textbf{PID Flashcards: Add Lab / Simulation Chapter Cards} \hfill \low \\
    Ch 10 (Simulation Lab) has 65 quiz questions and zero flashcards. Add 5--6 cards covering the Scilab step-test procedure, open-loop identification, and the tuning workflow from the lab.
\end{enumerate}

\smallskip
{\color{snavy}\rule{\linewidth}{0.4pt}}
\smallskip

{\small\sffamily\color{sgray}\MakeUppercase{Course Audit — SEL RTAC Gaps}}
\smallskip

\begin{enumerate}[resume]\small
  \item \textbf{RTAC Flashcards: Add Cards for 4 Uncovered Chapters} \hfill \high \\
    4 quiz chapters have zero flashcards: RTAC Architecture/Intro, ACSELERATOR Architect Software, Tag Database, and Modbus Gateway — the last one being syllabus Topic 4 with no flashcard coverage at all. Shoot for 5--7 cards per chapter; prioritize Gateway and TagDB since those are the heaviest operational topics.

  \item \textbf{RTAC: Add Lab Chapter Quiz Questions} \hfill \med \\
    The Lab \& Field Scenarios chapter (\texttt{Lab.jsx}) has a full web page and zero quiz questions — the only chapter across all 8 guides in that position. Write 30--40 questions covering the scenario walkthroughs on the lab page: configuration steps, fault identification, verification checks.

  \item \textbf{RTAC Flashcards: Reconcile Hardware Chapter Scope} \hfill \low \\
    The \texttt{hardware} flashcard category currently spans three web chapters (Architecture, Software, TagDB) instead of mapping to one. Refactor: move hardware-specific cards into an \texttt{intro} category and create separate \texttt{software} and \texttt{tagdb} categories so each chapter has its own flashcard home.
\end{enumerate}

\smallskip
{\color{snavy}\rule{\linewidth}{0.4pt}}
\smallskip

{\small\sffamily\color{sgray}\MakeUppercase{Course Audit — Ignition SCADA Gaps}}
\smallskip

\begin{enumerate}[resume]\small
  \item \textbf{Ignition Flashcards: Add Cards for 4 Uncovered Chapters} \hfill \high \\
    4 of 10 quiz chapters have zero flashcards: Intro/Overview (65 Qs), Ignition Gateway (65 Qs), Vision Module (65 Qs), and Alarming (65 Qs). That's 260 quiz questions with nothing backing them up. Priority order: (1) Gateway, (2) Alarming, (3) Vision, (4) Intro. Shoot for 7--8 cards per chapter.

  \item \textbf{Ignition Flashcards: Reconcile Quiz/Flashcard Naming} \hfill \low \\
    Quiz chapters use \texttt{scripting} and \texttt{alarms}; flashcard categories use \texttt{scripts} and \texttt{alarming}. A student going from a quiz result to the matching flashcard deck hits the wrong category. Standardize the naming across both data files so IDs line up.
\end{enumerate}

\smallskip
{\color{snavy}\rule{\linewidth}{0.4pt}}
\smallskip

{\small\sffamily\color{sgray}\MakeUppercase{Course Audit — Wireshark Gaps}}
\smallskip

\begin{enumerate}[resume]\small
  \item \textbf{Wireshark Flashcards: Add Advanced Techniques Chapter Cards} \hfill \med \\
    Ch 9 (Advanced Techniques) has 65 quiz questions and zero flashcards. The chapter covers editcap, mergecap, tshark, dumpcap ring buffer, Lua dissector scripting, and SSH remote capture. Add 6--8 cards hitting the most testable concepts: tshark command syntax, ring-buffer capture options, the SSH remote capture one-liner, and Lua dissector registration.

  \item \textbf{Wireshark Flashcards: Add Lab Chapter Cards} \hfill \low \\
    Ch 10 (Protocol Lab) has 65 quiz questions and zero flashcards. Add 5--6 cards covering the lab workflow: filter construction for Modbus, DNP3, and OPC UA captures; common frame-decode steps; and the ICS anomaly identification checklist from the lab exercises.

  \item \textbf{Wireshark Advanced Chapter: Add tcpreplay Coverage} \hfill \low \\
    \texttt{Advanced.jsx} covers editcap, mergecap, tshark, and SSH remote capture but says nothing about tcpreplay (PCAP replay for ICS protocol testing), which syllabus Topic 9 explicitly lists. Add a short section: how to replay a captured PCAP against a test device, rate control, and why it matters for testing IDS/IPS rules against replayed OT traffic.
\end{enumerate}

\smallskip
{\color{snavy}\rule{\linewidth}{0.4pt}}
\smallskip

{\small\sffamily\color{sgray}\MakeUppercase{Training Feature — Launch Checklist}}
\smallskip

\begin{enumerate}[resume]\small
  \item \textbf{Run Training Events SQL Migration} \hfill \high \\
    Open Supabase Dashboard $\to$ SQL Editor and run \texttt{scada-hub/supabase/20260516\_003\_training\_events.sql}. Creates the \texttt{training\_events} table, RLS policies, unique constraint on URL, and the course+date index. Has to happen before seeding or the \texttt{TrainingPanel} will break.

  \item \textbf{Configure Vercel Env Vars for Training Cron} \hfill \high \\
    In the scada-hub Vercel project settings, add these environment variables:\\
    \texttt{CRON\_SECRET} — any random string (e.g. \texttt{openssl rand -hex 32})\\
    \texttt{BRAVE\_SEARCH\_API\_KEY} — free tier at \texttt{api.search.brave.com} (2000 queries/month)\\
    \texttt{RESEND\_API\_KEY} — free account at \texttt{resend.com} (3000 emails/month)\\
    \texttt{ALERT\_EMAIL} — \texttt{ric3639@gmail.com} (broken URL alert destination)\\
    \texttt{SUPABASE\_SERVICE\_KEY} — service role key from Supabase project settings

  \item \textbf{Enable Vercel AI Gateway on scada-hub Project} \hfill \high \\
    In Vercel Dashboard $\to$ scada-hub project $\to$ Settings $\to$ AI Gateway: flip the gateway on. This provisions the OIDC token that \texttt{api/refresh-training.js} needs to call \texttt{anthropic/claude-haiku-4.5} without a direct API key in the environment.

  \item \textbf{Seed Initial Training Events} \hfill \high \\
    After the migration runs, seed 35 verified 2026 training events and certifications across all 8 courses:\\
    \texttt{SUPABASE\_URL=... SUPABASE\_SERVICE\_KEY=... node scripts/seed-training.js}\\
    Covers: Modbus (3), OPC UA (4), DNP3 (4), IEC 61131-3 (4), PID (5), RTAC (3), Ignition (5), Wireshark (6) + certs.

  \item \textbf{Deploy All 9 Projects — Training Feature} \hfill \high \\
    Redeploy all 9 projects to push the \texttt{TrainingPanel} component and Vercel cron config live. Run from each project root:\\
    \texttt{cd site \&\& GITHUB\_PAGES= npm run build \&\& cd .. \&\& vercel -{-}prod -{-}yes}\\
    scada-hub also needs: \texttt{vercel -{-}prod -{-}yes} from the project root (no site build needed for the cron function).

  \item \textbf{Troubleshooting Lesson Plan LaTeX PDF} \hfill \med \\
    Build a multi-chapter LaTeX PDF that turns real debugging sessions into field-engineer lesson plans. Each chapter covers one class of bug or architectural decision: symptom, what was checked, root cause, fix, and what to watch for next time. Example chapter: "Missing Icon Import Crashes React App" (Ignition \texttt{Zap} ReferenceError). Generate at the end of each session and stack chapters over time. Source lives in \texttt{scada-hub/lesson-plans/}. This is separate from the per-guide study guide PDFs — it's a \emph{developer training document}.
\end{enumerate}

\smallskip
{\color{snavy}\rule{\linewidth}{0.4pt}}
\smallskip

{\small\sffamily\color{sgray}\MakeUppercase{Professional Development / R\&D}}
\smallskip

\begin{enumerate}[resume]\small
  \item \textbf{ElevenLabs TTS Audio Regeneration} \hfill \med \\
    \texttt{normalizeForTTS()} function written and wired into \texttt{generate-audio-all.mjs} — symbols now spoken as words (\texttt{=} $\to$ equals, \texttt{!=} $\to$ does not equal, \texttt{0x03} $\to$ hex 03, \texttt{:4840} $\to$ port 4840, etc.). \textbf{Remaining:} Regenerate 178 modbus MP3s with normalized text. ElevenLabs quota was at zero when regen was attempted (2026-05-18) — backup preserved at \texttt{site/public/audio-backup/}. Run \texttt{node generate-audio-all.mjs modbus-study-guide} from \texttt{\textasciitilde/Downloads/} when monthly quota resets. Other 7 guides have no audio yet.

  \item \textbf{Badge/Achievement System} \hfill \done \\
    Shipped 2026-05-18. Six flat-top hexagonal badges per guide: \textit{First Spark} (amber, any L1 pass), \textit{Foundations Laid} (blue, all L1), \textit{Field Technician} (orange, all L2), \textit{Scholar} (purple, any L3), \textit{War Stories} (green, any Field Scenario), and a guide-specific certified badge (gold, all L1+L2+L3). Silhouettes shown for unearned; earned badges glow with per-badge color. Animated modal overlay (spring entrance, breathing glow, flavor text) fires once per badge on first earn. Keyboard-dismissible (Esc, Enter, Space). All 8 guides committed and pushed.

  \item \textbf{Wireshark Blank Page (Vercel)} \hfill \done \\
    Production crash: \texttt{ReferenceError: Can't find variable: BookOpen} in \texttt{Sidebar.jsx}. Icon was used in the Chapters dropdown button but omitted from the lucide-react import. Dev server resolved it silently; production bundle did not. Added \texttt{ErrorBoundary} to surface the error, fixed the import, removed the boundary. Shipped 2026-05-19.

  \item \textbf{All Courses Cross-Guide Links} \hfill \done \\
    Removed environment-conditional \texttt{isGH} logic from all 8 sidebars. Links now always point to \texttt{\$\{slug\}-study-guide.vercel.app} — all 8 guides have confirmed Vercel deployments. Shipped 2026-05-19.

  \item \textbf{OPC UA Orphaned GIF --- Intro.jsx} \hfill \done \\
    \texttt{courseHero} GifCard rendered as \texttt{flex justify-end} with empty left half. Wrapped in \texttt{flex items-start gap-6} div with body paragraph. Shipped 2026-05-19.

  \item \textbf{Unified Streak Counter} \hfill \done \\
    Per-guide streak keys (\texttt{wireshark\_streak}, \texttt{rtac\_streak}, etc.) replaced with shared \texttt{scadahub\_streak} / \texttt{scadahub\_streak\_date} across all 8 \texttt{Sidebar.jsx} and \texttt{Home.jsx} files. One streak advances across any guide visit. Shipped 2026-05-19.

  \item \textbf{Landing Page Redesign --- More Professional} \hfill \med \\
    The current scada-hub landing page has been described as looking like an ``Instagram ad.'' Needs a redesign toward a professional training platform aesthetic: cleaner layout, less promotional language, tighter typography, less heavy-handed animation. Reference: GitBook, Udemy Business, Pluralsight.

  \item \textbf{R\&D: AI Agent Skills and Agentic Tooling} \hfill \low \\
    Get real reps in with agentic AI tools that actually matter for this stack — specifically Claude AI (Claude Code, agent SDK, multi-agent orchestration), OpenAI Codex, and equivalent code-generation toolchains. Figure out which capabilities can move the needle on delivery (code generation, automated testing, documentation, diagram synthesis) and build repeatable workflows around the ones that stick. Deliverable: a short internal write-up covering what was evaluated, what it can actually do in this stack, and what's worth adopting.
\end{enumerate}

\end{document}
